\section*{Propuesta Indice}

\setlist{nosep, leftmargin=*} 

\begin{enumerate}
  % --- CAPÍTULO 1 ---
  \item \textbf{Introducción general}
  \begin{enumerate}[label*=\arabic*.]
    \item Contexto y motivación del estudio
    \item Problemática de los efectos de escala en la predicción de potencia
    \item Importancia para los buques pesqueros y la industria nacional (?)
    \item Objetivos generales y específicos
    \item Estructura de la tesis
  \end{enumerate}

  % --- CAPÍTULO 2 (Aquí estaba el error) ---
  \item \textbf{Fundamentos teóricos y Estado del Arte}
  \begin{enumerate}[label*=\arabic*.] 
    \item Fundamentos de Mecánica de Fluidos 
    \item Teoría de la Capa Límite y Turbulencia 
    \item Componentes de la Resistencia al Avance 
    \item Teoría de Semejanza y Efectos de Escala 
    \item Métodos Numéricos en Hidrodinámica (CFD) 
    \item Aplicaciones combinadas CFD/EFD en predicción de potencia 
  \end{enumerate} 

  % --- CAPÍTULO 3 ---
  \item \textbf{Metodología general}
  \begin{enumerate}[label*=\arabic*.]
    \item Enfoque combinado CFD/EFD
    \item Instalaciones experimentales (LabHiNO)
    \item Método de Volúmenes Finitos y software utilizado (OpenFOAM)
    \item Estrategia de verificación y validación (V\&V)
    \item Descripción de los buques y hélices de referencia
    \item Descripción de los buques pesqueros analizados
    \item Estructura metodológica por ejes
  \end{enumerate}

  % --- CAPÍTULO 4 ---
  \item \textbf{Eje 1: Resistencia al avance — Factor de forma}
  \begin{enumerate}[label*=\arabic*.]
    \item Determinación experimental y numérica del factor de forma
    \begin{enumerate}[label*=\arabic*.]
      \item Ensayos experimentales de resistencia al avance 
      \item Determinación experimental del factor de forma 
      \item Ensayos computacionales de resistencia al avance
      \item Determinación computacional del factor de forma
      \item Comparación EFD–CFD de coeficientes de resistencia y de factor de forma
      \item Discusión de incertidumbres y diferencias entre EFD y CFD
      \item Análisis de efectos de escala
      \item Conclusiones parciales
    \end{enumerate}

    \item Influencia de la proa bulbo en el factor de forma
    \begin{enumerate}[label*=\arabic*.]
      \item Motivación y antecedentes sobre proa bulbo en pesqueros
      \item Modelos numéricos 
      \item Resultados de resistencia para diferentes configuraciones de proa
      \item Efecto sobre el factor de forma
      \item Conclusiones parciales
    \end{enumerate}
    
    \item Influencia de la popa sumergida en el factor de forma
    \begin{enumerate}[label*=\arabic*.]
       \item Introducción y motivación del estudio
       \item Definiciones geométricas y parámetros característicos
       \item Desarrollo del método 2-k
       \item Aplicación al pesquero de referencia
       \item Análisis comparativo de flujo
       \item Conclusiones parciales
    \end{enumerate}    
    
    \item Influencia del modelo de turbulencia en la determinación del factor de forma
    \begin{enumerate}[label*=\arabic*.]
       \item Introducción y motivación del estudio
       \item Metodología y configuración numérica
       \item Modelos de turbulencia analizados ($k$–$\varepsilon$, $k$–$\varepsilon$ realizable, $k$–$\omega$, $k$–$\omega$ SST)
       \item Resultados comparativos de resistencia viscosa
       \item Efecto sobre el factor de forma
       \item Descomposición de resistencia según origen friccional y de presión viscosa
       \item Recomendaciones para selección de modelos en estudios de escala
       \item Conclusiones parciales
    \end{enumerate}

    \item Efecto de la relación L/B en el factor de forma
    \begin{enumerate}[label*=\arabic*.]
      \item Motivación y observaciones sobre L/B en cascos de baja esbeltez
      \item Presentación de pesqueros con distintos L/B
      \item Análisis CFD y comparación de distribución de presión y resistencia viscosa
      \item Generalización del comportamiento observado y formulación empírica
      \item Implicancias para el escalado y la extrapolación de resultados
      \item Conclusiones parciales
    \end{enumerate}
  \end{enumerate}

  % --- CAPÍTULO 5 ---
  \item \textbf{Eje 2: Propulsión — Ensayos de hélices en aguas abiertas}
  \begin{enumerate}[label*=\arabic*.]
    \item Metodología para el estudio de hélices en condiciones de aguas abiertas
    \item Resultados experimentales de hélices en ensayos de aguas abiertas
    \item Resultados computacionales de hélices en ensayos de aguas abiertas
    \item Comparación de métodos de extrapolación ITTC con CFD en escala real
    \item Efectos de escala  
    \item Efectos del modelo de turbulencia (turbulento o de transición) 
    \item Efecto de la rugosidad y del agregado de turbuladores
    \item Conclusiones parciales 
  \end{enumerate}

  % --- CAPÍTULO 6 ---
  \item \textbf{Síntesis y aportes finales}
  \begin{enumerate}[label*=\arabic*.]
    \item Integración de resultados de resistencia y propulsión
    \item Evaluación global de los métodos combinados CFD/EFD
    \item Aportes sobre efectos de escala en pesqueros
    \item Implicancias para la práctica de diseño y ensayos
    \item Líneas futuras de investigación
  \end{enumerate}

  \item \textbf{Conclusiones generales}

  \item \textbf{Referencias bibliográficas}

  \item \textbf{Apéndices}
  \begin{enumerate}[label*=\Alph*.]
    \item Detalles experimentales de los ensayos en el LabHiNO
    \item Condiciones numéricas y mallado de los casos CFD
    \item Tablas de resultados adicionales
    \item Publicaciones
  \end{enumerate}
\end{enumerate}